\documentclass[a4paper]{article}

% \usepackage[default]{fontsetup}

\usepackage{fancyhdr}
\usepackage{extramarks}
\usepackage{amsmath}
\usepackage{amsthm}
\usepackage{amsfonts}
\usepackage{tikz}
\usepackage[plain]{algorithm}
\usepackage{algpseudocode}
\usepackage{enumerate}
\usepackage{tikz}

\usetikzlibrary{automata,positioning}

%
% Basic Document Settings
%  

\topmargin=-0.2in
\evensidemargin=0in
\oddsidemargin=0in
\textwidth=6.5in
\textheight=9.5in
\headsep=0.25in

\linespread{1.1}

\pagestyle{fancy}
\lhead{\hmwkAuthorName}
\chead{\hmwkClass : \hmwkTitle}
\rhead{\firstxmark}
\lfoot{\lastxmark}
\cfoot{\thepage}

\renewcommand\headrulewidth{0.4pt}
\renewcommand\footrulewidth{0.4pt}

\setlength\parindent{0pt}

%
% Create Problem Sections
%

\newcommand{\enterProblemHeader}[1]{
    \nobreak\extramarks{}{Problem \arabic{#1} continued on next page\ldots}\nobreak{}
    \nobreak\extramarks{Problem \arabic{#1} (continued)}{Problem \arabic{#1} continued on next page\ldots}\nobreak{}
}

\newcommand{\exitProblemHeader}[1]{
    \nobreak\extramarks{Problem \arabic{#1} (continued)}{Problem \arabic{#1} continued on next page\ldots}\nobreak{}
    \stepcounter{#1}
    \nobreak\extramarks{Problem \arabic{#1}}{}\nobreak{}
}

\newcommand*\circled[1]{\tikz[baseline=(char.base)]{
		\node[shape=circle,draw,inner sep=2pt] (char) {#1};}}


\setcounter{secnumdepth}{0}
\newcounter{partCounter}
\newcounter{homeworkProblemCounter}
\setcounter{homeworkProblemCounter}{1}
\nobreak\extramarks{Problem \arabic{homeworkProblemCounter}}{}\nobreak{}

%
% Homework Problem Environment
%
% This environment takes an optional argument. When given, it will adjust the
% problem counter. This is useful for when the problems given for your
% assignment aren't sequential. See the last 3 problems of this template for an
% example.
%

\newenvironment{homeworkProblem}[2][-1]{%
    \ifnum#1>0\relax
        \setcounter{homeworkProblemCounter}{#1}
    \fi
    \section{Problem \arabic{homeworkProblemCounter} (#2 points)}
    \setcounter{partCounter}{1}
    \enterProblemHeader{homeworkProblemCounter}
}{%
    \exitProblemHeader{homeworkProblemCounter}
}

%
% Homework Details
%   - Title
%   - Class
%   - Due date
%   - Name
%   - Student ID

\newcommand{\hmwkTitle}{Homework\ \#05}
\newcommand{\hmwkClass}{Probability \& Statistics for EECS}
\newcommand{\hmwkDueDate}{2026-10-25}
\newcommand{\hmwkAuthorName}{Zhang Yiheng}
\newcommand{\hmwkAuthorID}{2025531068}


%
% Title Page
%

\title{
    \vspace{2in}
    \textmd{\textbf{\hmwkClass:\\  \hmwkTitle}}\\
    \normalsize\vspace{0.1in}\small{Due\ on\ \hmwkDueDate\ at 23:59}\\
	\vspace{4in}
}

\author{
	Name: \textbf{\hmwkAuthorName} \\
	Student ID: \hmwkAuthorID}
\date{}

\renewcommand{\part}[1]{\textbf{\large Part \Alph{partCounter}}\stepcounter{partCounter}\\}

%
% Various Helper Commands
%

% Useful for algorithms
\newcommand{\alg}[1]{\textsc{\bfseries \footnotesize #1}}
% For derivatives
\newcommand{\deriv}[1]{\frac{\mathrm{d}}{\mathrm{d}x} (#1)}
% For partial derivatives
\newcommand{\pderiv}[2]{\frac{\partial}{\partial #1} (#2)}
% Integral dx
\newcommand{\dx}{\mathrm{d}x}
% Alias for the Solution section header
\newcommand{\solution}{\textbf{\large Solution}}
% Probability commands: Expectation, Variance, Covariance, Bias
\newcommand{\E}{\mathrm{E}}
\newcommand{\Var}{\mathrm{Var}}
\newcommand{\Cov}{\mathrm{Cov}}
\newcommand{\Bias}{\mathrm{Bias}}

\begin{document}


% \maketitle
% \thispagestyle{empty}
% \pagebreak

\date{
Due on Oct. 25, 2026, 23:59 UTC+8}
\title{SI 140A.02  Probability \& Statistics for EECS, Fall 2026 \\
Homework 5}
\maketitle
Read all the instructions below carefully before you start working on the assignment, and before you make a submission.
\begin{itemize}
    \item You are required to submit your homework in PDF format generated by \textbf{LaTeX}, otherwise you will get no points of this homework.
    \item You are required to write down all the major steps towards making your conclusions; otherwise you may obtain limited points of the problem.
    \item Write your homework in English; otherwise you will get no points of this homework.
    \item Any form of plagiarism will lead to $0$ point of this homework. 
\end{itemize}
\newpage

\begin{homeworkProblem}{20}

A/B testing is a randomized experiment used to compare two
treatments, Treatment A and Treatment B. Users arrive one by one
and are assigned to one of the treatments. Each trial is classified
as either a success or a failure.

Consider the following ``stay-with-a-winner'' procedure:
\begin{enumerate}[(i)]
    \item Assign the first user to Treatment A or Treatment B
    with equal probabilities.
    \item If the trial for the $n$th user is a success, assign
    the $(n+1)$st user to the same treatment. Otherwise, switch
    to the other treatment.
\end{enumerate}

The success probabilities for Treatments A and B are $a$ and $b$,
respectively, where $a\ne b$. Given the assigned treatment, these
success probabilities do not depend on the previous outcomes.

Let $p_n$ be the probability of success on the $n$th trial, and
let $a_n$ be the probability that Treatment A is assigned on the
$n$th trial.

\begin{enumerate}[(a)]
    \item Show that
    \[
        p_n=(a-b)a_n+b,
        \qquad
        a_{n+1}=(a+b-1)a_n+1-b.
    \]

    \item Use part (a) to show that
    \[
        p_{n+1}=(a+b-1)p_n+a+b-2ab.
    \]

    \item Assuming that the limit exists, find
    \[
        \lim_{n\to\infty}p_n.
    \]
\end{enumerate}

\solution

\part{a}
Conditioning on the treatment assigned on trial $n$,
\[
    p_n=a\,a_n+b(1-a_n)=(a-b)a_n+b.
\]
For $a_{n+1}$, note that Treatment A is used on trial $n+1$ if
either A was used on trial $n$ and succeeded, or B was used on
trial $n$ and failed. Hence
\[
    a_{n+1}=a\,a_n+(1-b)(1-a_n)=(a+b-1)a_n+1-b.
\]

\part{b}
Using part (a) twice,
\begin{align*}
    p_{n+1}
    &=(a-b)a_{n+1}+b
     =(a-b)\bigl[(a+b-1)a_n+1-b\bigr]+b\\
    &=(a+b-1)(p_n-b)+(a-b)(1-b)+b\\
    &=(a+b-1)p_n-b(a+b-1)+(a-b)(1-b)+b\\
    &=(a+b-1)p_n+a+b-2ab.
\end{align*}

\part{c}
If $p_n\to p$, taking limits in part (b) gives
\[
    p=(a+b-1)p+a+b-2ab
    \;\Longrightarrow\;
    (2-a-b)p=a+b-2ab,
\]
and therefore
\[
    \lim_{n\to\infty}p_n=\frac{a+b-2ab}{2-a-b}.
\]

\end{homeworkProblem}

\newpage

\begin{homeworkProblem}{20}

A fair coin is tossed repeatedly and independently, producing a
sequence of heads and tails. Let $N$ be the number of tosses
required to observe the first occurrence of the pattern $HH$.

Find the PMF of $N$.

\solution

Let $f_n=\Pr(N=n)$. Clearly $f_1=0$ and $f_2=\Pr(HH)=1/4$.

For $n\ge3$, condition on the first two tosses:
\begin{itemize}
    \item if the first toss is $T$ (probability $1/2$), the process
    restarts, contributing $\frac12 f_{n-1}$;
    \item if the first two tosses are $HT$ (probability $1/4$), the
    process restarts after two tosses, contributing
    $\frac14 f_{n-2}$.
\end{itemize}
Hence
\[
    f_n=\frac12 f_{n-1}+\frac14 f_{n-2},\qquad n\ge3.
\]
The characteristic equation $r^2=\frac12 r+\frac14$ has roots
\[
    r=\frac{1\pm\sqrt5}{4}.
\]
Writing $\varphi=\frac{1+\sqrt5}{2}$ and
$\psi=\frac{1-\sqrt5}{2}$, the general solution is
\[
    f_n=A\left(\frac{\varphi}{2}\right)^n
        +B\left(\frac{\psi}{2}\right)^n.
\]
The initial conditions $f_1=0$ and $f_2=1/4$ give
\[
    A\varphi+B\psi=0,
    \qquad
    A\varphi^2+B\psi^2=1,
\]
whose solution is
\[
    A=\frac{1}{\sqrt5\,\varphi},
    \qquad
    B=-\frac{1}{\sqrt5\,\psi}.
\]
Therefore
\begin{align*}
    f_n
    &=\frac{1}{\sqrt5}\left(\frac{\varphi^{n-1}}{2^n}
      -\frac{\psi^{n-1}}{2^n}\right)
    =\frac{\varphi^{n-1}-\psi^{n-1}}{\sqrt5\,2^n}
    =\frac{F_{n-1}}{2^n},
\end{align*}
where
\[
    F_k=\frac{\varphi^k-\psi^k}{\sqrt5}
\]
is the $k$th Fibonacci number ($F_0=0$, $F_1=F_2=1$). Hence
\[
    \Pr(N=n)=\frac{F_{n-1}}{2^n},\qquad n\ge1.
\]

\end{homeworkProblem}

\newpage


\begin{homeworkProblem}{15}

A system contains seven devices: four in Group A and three in
Group B. Each device fails independently with probability $p$,
where $0<p<1$.

Let $X$ be the number of failed devices in Group A, and let $T$
be the total number of failed devices in the system.

\begin{enumerate}[(a)]
    \item Find the PMF of $X$ and the PMF of $T$.

    \item Given that exactly three devices have failed, find
    the conditional PMF of $X$. Identify the distribution.

    \item Explain why the conditional distribution in part (b)
    does not depend on $p$.
\end{enumerate}

\solution

\part{a}
Each device fails independently with probability $p$, so
\[
    \Pr(X=i)=\binom{4}{i}p^i(1-p)^{4-i},\qquad i=0,1,2,3,4,
\]
and
\[
    \Pr(T=j)=\binom{7}{j}p^j(1-p)^{7-j},\qquad j=0,1,\dots,7.
\]
That is, $X\sim\mathrm{Binomial}(4,p)$ and
$T\sim\mathrm{Binomial}(7,p)$.

\part{b}
For $i=0,1,2,3$,
\begin{align*}
    \Pr(X=i\mid T=3)
    &=\frac{\Pr(T=3\mid X=i)\Pr(X=i)}{\Pr(T=3)}\\
    &=\frac{\binom{3}{3-i}p^{3-i}(1-p)^{i}
            \binom{4}{i}p^{i}(1-p)^{4-i}}
           {\binom{7}{3}p^{3}(1-p)^{4}}\\
    &=\frac{\binom{4}{i}\binom{3}{3-i}}{\binom{7}{3}},
\end{align*}
where $\Pr(T=3\mid X=i)$ corresponds to exactly $3-i$ failures
among the three devices of Group B. Numerically,
\[
    \Pr(X=i\mid T=3)
    =\frac{1}{35},\ \frac{12}{35},\ \frac{18}{35},\ \frac{4}{35}
    \qquad (i=0,1,2,3).
\]
This is the hypergeometric distribution: given that three of the
seven devices failed, the failing devices form a uniformly random
three-element subset of the seven, and $X$ counts how many of
them belong to the four devices of Group A.

\part{c}
Since all seven devices have the same failure probability $p$
and fail independently, conditioning on the total number of
failures makes every subset of that size equally likely; the
value of $p$ affects only how many devices fail, not which ones
fail. Thus the conditional PMF in part (b) does not involve $p$.

\end{homeworkProblem}

\newpage

\begin{homeworkProblem}{15}

Two fair six-sided dice are rolled independently. Let $U$ and
$V$ be their respective outcomes, and define
\[
    D=|U-V|.
\]

\begin{enumerate}[(a)]
    \item Find the PMF of $D$.

    \item Find the CDF $F_D(t)$ for every real number $t$.

    \item Define $W=\min(D,2)$. Find the PMF of $W$.
\end{enumerate}

\solution

\part{a}
There are $36$ equally likely outcomes $(U,V)$. For $i=1,\dots,5$
there are $2(6-i)$ outcomes with $|U-V|=i$, while $D=0$ occurs
for the $6$ outcomes with $U=V$. Hence
\[
    \Pr(D=i)=
    \begin{cases}
        \dfrac{1}{6}, & i=0,\\[1mm]
        \dfrac{6-i}{18}, & i=1,2,3,4,5,\\[1mm]
        0, & \text{otherwise}.
    \end{cases}
\]

\part{b}
Accumulating the PMF from part (a),
\[
    F_D(t)=
    \begin{cases}
        0, & t<0,\\[1mm]
        \dfrac{1}{6}, & 0\le t<1,\\[1mm]
        \dfrac{4}{9}, & 1\le t<2,\\[1mm]
        \dfrac{2}{3}, & 2\le t<3,\\[1mm]
        \dfrac{5}{6}, & 3\le t<4,\\[1mm]
        \dfrac{17}{18}, & 4\le t<5,\\[1mm]
        1, & t\ge5.
    \end{cases}
\]

\part{c}
Since $W=\min(D,2)$,
\[
    \Pr(W=0)=\Pr(D=0)=\frac16,\qquad
    \Pr(W=1)=\Pr(D=1)=\frac{5}{18},
\]
and
\[
    \Pr(W=2)=\Pr(D\ge2)=1-\frac16-\frac{5}{18}=\frac{5}{9}.
\]
Thus $\Pr(W=w)=0$ for $w\notin\{0,1,2\}$.

\end{homeworkProblem}

\newpage

\begin{homeworkProblem}{15}

There are two coins, A and B. Coin A lands heads with probability
$1/4$, while coin B lands heads with probability $3/4$.

One coin is selected with equal probabilities and then tossed
twice. The same selected coin is used for both tosses. Given
the selected coin, the two tosses are independent.

Let $C\in\{A,B\}$ denote the selected coin. Let $X$ and $Y$ be
the indicator random variables for heads on the first and
second tosses, respectively.

\begin{enumerate}[(a)]
    \item Find the PMF of $X$ and the PMF of $Y$.

    \item Are $X$ and $Y$ independent? Are they conditionally
    independent given $C$? Justify your answers.

    \item Given that both tosses are heads, find the
    conditional probability that coin B was selected.
\end{enumerate}

\solution

\part{a}
By the law of total probability,
\[
    \Pr(X=1)=\frac12\cdot\frac14+\frac12\cdot\frac34=\frac12,
\]
so $\Pr(X=0)=\frac12$. The two tosses use the same coin and
therefore have the same marginal distribution, so $Y$ has the
same PMF:
\[
    \Pr(Y=1)=\Pr(Y=0)=\frac12.
\]

\part{b}
The variables $X$ and $Y$ are not independent. Indeed,
\[
    \Pr(X=1,Y=1)
    =\frac12\left(\frac14\right)^2
     +\frac12\left(\frac34\right)^2
    =\frac{1}{32}+\frac{9}{32}
    =\frac{5}{16},
\]
whereas
\[
    \Pr(X=1)\Pr(Y=1)=\frac12\cdot\frac12=\frac14\ne\frac{5}{16}.
\]
They are conditionally independent given $C$, because given the
selected coin the two tosses are independent:
\[
    \Pr(X=x,Y=y\mid C)=\Pr(X=x\mid C)\Pr(Y=y\mid C)
\]
for all $x,y\in\{0,1\}$.

\part{c}
By Bayes' rule,
\begin{align*}
    \Pr(C=B\mid X=1,Y=1)
    &=\frac{\Pr(X=1,Y=1\mid C=B)\Pr(C=B)}
           {\Pr(X=1,Y=1)}\\
    &=\frac{\frac{9}{16}\cdot\frac12}{\frac{5}{16}}
    =\frac{9}{10}.
\end{align*}

\end{homeworkProblem}

\newpage


\begin{homeworkProblem}{15}

A random variable $X$ has PMF
\[
    \Pr(X=1)=\frac12,\qquad
    \Pr(X=2)=\frac14,\qquad
    \Pr(X=3)=\Pr(X=4)=\frac18.
\]
Let $U$ have the discrete uniform distribution on
$\{1,2,3,4\}$.

Throughout this problem, logarithms have base 2, and entropy
is measured in bits.

\begin{enumerate}[(a)]
    \item Find $H(X)$ and $H(U)$.

    \item Define
    \[
        Y=
        \begin{cases}
            1,&X\le2,\\
            0,&X>2,
        \end{cases}
        \qquad
        Z=
        \begin{cases}
            1,&X=1,\\
            0,&X\ne1.
        \end{cases}
    \]
    Find the PMF of each variable and calculate $H(Y)$
    and $H(Z)$.

    \item Is the size of the support sufficient to determine
    the entropy of a discrete random variable?
    Explain using distributions from this problem.
\end{enumerate}

\solution

\part{a}
\[
    H(X)=-\sum_{i=1}^{4}\Pr(X=i)\log_2\Pr(X=i)
    =\frac12\log_2 2+\frac14\log_2 4
     +2\cdot\frac18\log_2 8
    =\frac12+\frac12+\frac34=\frac74=1.75\text{ bits},
\]
and since $U$ is uniform on a set of size $4$,
\[
    H(U)=\log_2 4=2\text{ bits}.
\]

\part{b}
Since $\Pr(Y=1)=\Pr(X\le2)=\frac12+\frac14=\frac34$ and
$\Pr(Y=0)=\frac14$,
\[
    H(Y)
    =-\frac34\log_2\frac34-\frac14\log_2\frac14
    =\frac34\log_2\frac43+\frac12
    =2-\frac34\log_2 3
    \approx0.811\text{ bits}.
\]
Since $\Pr(Z=1)=\Pr(X=1)=\frac12$ and $\Pr(Z=0)=\frac12$,
\[
    H(Z)=-\frac12\log_2\frac12-\frac12\log_2\frac12
    =1\text{ bit}.
\]

\part{c}
No. The support size only gives the upper bound
$H\le\log_2|\text{support}|$, with equality exactly for the
uniform distribution on the support; it does not determine the
entropy. In this problem $X$ and $U$ both have support
$\{1,2,3,4\}$, but $H(X)=1.75<2=H(U)$. Similarly, $Y$ and $Z$
both have a support of size $2$, but
$H(Y)\approx0.811<1=H(Z)$.

\end{homeworkProblem}

\end{document}